\section{Almost reflective dependently typed language:
    MLTT-R}
\label{MLTT-R}

Now we are ready to take on the main challenge: design a
minimal dependently typed reflective language, and
provide internal models of itself to yield universe
polymorphism, external parametricity and do synthetic
mathematics inside it. As a starting point, we take
Martin-L\"of type theory with $\Pi-$,$\Sigma-$,
$\mathrm{Id}-$ types, natural number type and infinite
non-cumulative hierarchy of universes. As discussed in
section \ref{no-go-section}, the best practical solution
to reinterpretation here would be to forbid reductions
under it and reinterpret terms as-is, and constrain
models using $\mathrm{Id}$-types. Still, in comparison
to non-dependently typed case, there are infinitely more
problems to solve if we want to define a truly
reflective system:

\begin{enumerate}
    \item If we na\"{\i}vely adopt an approach to typing
        of models from section \ref{HM-R}, in order to
        specify a model, we would need to specify
        countably many models, one per each universe.
        And even if we did, to which universe does the
        type of models belong?
    \item Now that we have computations in types, we
        have to deal with conversion rule under
        reinterpretation:
        \[
            \dfrac{t:T\quad T\equiv T'}{t:T'}
            \stackrel{?}{\leadsto}
            \dfrac{\llbracket t\rrbracket:\llbracket T
            \rrbracket\quad\llbracket T\rrbracket\equiv
            \llbracket T'\rrbracket}{\llbracket t
            \rrbracket:\llbracket T'\rrbracket}
        \]
        The problem is that, after interpretation,
        $\llbracket T \rrbracket$ is only
        \textit{propositionally} equal to $\llbracket T'
        \rrbracket$.
    \item If we want a truly reflective system, we need
        to know what to do with \textit{reinterpretation
        of reinterpretation}:
        \[
            \lambda x.\llbracket\llbracket x\rrbracket_1
            \rrbracket_2
            \quad\text{and}\quad
            \llbracket\lambda x.\llbracket x\rrbracket_1
            \rrbracket_2
        \]
        In the first case, we perform two subsequent
        interpretations on a yet unknown $x$; in the
        second case, we interpret the process of
        interpretation itself. What should this mean?
\end{enumerate}

To deal with universe sizing issues, consider the
following observation: in a system without level
binders, each term uses only a finite number of levels.
Hence to reinterpret this term, we need to only provide
a finite number of universe models, which can together
fit into some larger universe of a finite level. No
universe polymorphism required!

To interpret conversion rule, we start by first
pinpointing the locations where conversion rule is
actually used; in order to do this, we \textit{need} to
use bidirectional type checking. At the moment of
conversion, interpretation has enough information to
know both $T$ and $T'$; thus we can translate a chain of
reductions performed during definitional equality check
into a chain of propositional equalities taken from
coherence conditions proven in the definition of a
model. Then, a point of conversion is wrapped in
equality elimination with synthesized witness as a
scrutinee.

In reinterpreting reinterpretation, we actually only
deal with ``double reinterpretation'': cases like
$\lambda x.\llbracket\llbracket x\rrbracket_1
\rrbracket_2$ can be dealt with by \textit{merging} two
models into a new one, a technique reminiscent of
\textit{Collapsing towers of interpreters} (TODO: cite).
We do not try to actually interpret interpretation
itself as in $\llbracket\lambda x.\llbracket x
\rrbracket_1\rrbracket_2$ and introduce a modality
$\Box$ meaning ``no interpretation was used in
construction of this term''; reinterpretable terms have
to have a boxed type. We also use $\Box$ to relax a ban
on computations inside reinterpretation by allowing
``splicing'' boxed expressions into other boxed
expressions; inside splices, computation is allowed to
bring terms to the boxed form for further
interpretation.

Note that discussion above was ambient to the choice of
formalism to present models. In line with current state
of the art, we use Categories-with-Families
\cite{castellan2020cwf} as one of the most convenient
frameworks used today to specify models of dependent
type theories.

TODO: CwF definition

To actually compute interpretation in MLTT-R, we need
two mutually recursive functions:
$t\,@^\gamma_{\Leftarrow T}\mathcal{U}$ and
$t\,@^\gamma_{\Rightarrow}\mathcal{U}$ for checkable and
synthesizing terms, respectively. TODO: FINISH.

\begin{figure}
\centering
\begin{bnf}
$\Gamma$ : Context ::=
    | $\varepsilon$ : empty global context
    | $\Gamma;\Theta$ : localization
    ;;
$\Theta$ : Small context ::=
    | $\cdot$ : empty small context
    | $\Theta$, $x\,\colon T$ : variable extension
    ;;
$T$, $S$ : Type ::=
    | $\mathcal{U}_\ell$ : universe
    | $\mathrm{El}_\ell(t)$ : universal type family
    | $\mathrm{Nat}_\ell$ : natural number type
    | $(x\,\colon S)\times T$ : $\Sigma$-type
    | $(x\,\colon S)\to T$ : $\Pi$-type
    | $t_1=_T t_2$ : identity type
    | $\Box\,T$ : type of closed terms
    | $\mathcal{M}_{\ell_0,\ldots,\ell_\kappa}$
        : type of models
    | $T^t_\Gamma$ : typal reinterpretation
    ;;
$\ell$, $\kappa$, $j$ : Level :in: $\mathbb{N}$
    : finite levels ;;
$t$, $s$ : Term ::=
    | $x$ // $y$ // $z$ // \texttt{m} // $\ldots$
        : variables
    | $\mathtt{U}_\ell$ : code of a universe
    | $\mathtt{Nat}_\ell$
        : code of a natural number type
    | $0$ : zero
    | $\mathrm{S}\,t$ : successor
    | $\mathrm{natind}_{x.T}\,t\,\{s_0,y\mapsto s_S\}$
        : induction
    | $\mathtt{Sigma}(x\,\colon s)t$
        : code of a $\Sigma$-type
    | $(s, t)$ : pairing
    | $\pi_1\,t$ // $\pi_2\,t$ : projections
    | $\mathtt{Pi}(x\,\colon s)t$
        : code of a $\Pi$-type
    | $\lambda x.\,t$ : abstraction
    | $t\,s$ : application
    | $\mathtt{Id}\,t\,t_1\,t_2$
        : code of an identity type
    | $\mathrm{refl}$ : reflexivity
    | $\mathrm{J}_{xy.S}\,t\,s$
        : J rule (elimination of equality type)
    | $\mathtt{Box}\,t$ : code of a box type
    | $[t]_\text{\faLock}$ : locking
    | $[t]_\text{\faUnlock}$ : unlocking
    | $\llbracket t
        \rrbracket^{\ell_0,\ldots,\ell_\kappa}$
        : internal interpretation
    | $t_\Gamma^s$ : intermediate interpretation
    | $t\,\colon T$ : ascription
    ;;
\end{bnf}
\caption{MLTT-R grammar}
\Description{}
\label{mltt-r-grammar}
\end{figure}

\begin{figure}
\centering
\begin{mathpar}
    \inferrule[WF-Loc]{\Gamma\vdash\Theta}
        {\Gamma;\Theta\;\mathrm{wf}}
    \and
    \inferrule[WF-Eps]{}{\varepsilon\vdash\cdot}
    \and
    \inferrule[WF-Nil]{\Gamma\;\mathrm{wf}}
        {\Gamma\vdash\cdot}
    \and
    \inferrule[WF-Var]{
        \Gamma\vdash\Theta\\
        \Gamma;\Theta\vdash\mathrm{level}(T)=\ell
    }{\Gamma\vdash\Theta,x:T}
    \\
    \inferrule[Level-Uni]{\Gamma\;\mathrm{wf}}{
        \Gamma\vdash\mathrm{level}(
            \mathcal{U}_\ell) = \ell+1
    }
    \and
    \inferrule[Level-El]
        {\Gamma\vdash t\Leftarrow\mathcal{U}_\ell}
        {\Gamma\vdash\mathrm{level}(
            \mathrm{El}_\ell(t)) = \ell}
    \and
    \inferrule[Level-Nat]{\Gamma\;\mathrm{wf}}{
        \Gamma\vdash\mathrm{level}(
            \mathrm{Nat}_\ell) = \ell}
    \and
    \inferrule[Level-Box]
        {\Gamma;\cdot\vdash\mathrm{level}(T)=\ell}
        {\Gamma\vdash\mathrm{level}(\Box\,T)=\ell}
    \and
    \inferrule[Level-Sigma, Level-Pi]{
        \Gamma;\Theta\vdash\mathrm{level}(S) = \kappa \\
        \Gamma;\Theta,x:S\vdash\mathrm{level}(T) = \ell
    }{
        \Gamma;\Theta\vdash
            \mathrm{level}((x:S)\times T)=
            \mathrm{level}((x:S)\to T)=
            \max(\kappa,\ell)
    }
    \and
    \inferrule[Level-Id]{
        \Gamma\vdash\mathrm{level}(T)=\ell \\
        \Gamma\vdash t_1,t_2\Leftarrow T
    }{\Gamma\vdash\mathrm{level}(t_1=_T t_2)=\ell}
    \and
    \inferrule[Level-Model]{\Gamma\;\mathrm{wf}}{
        \Gamma\vdash\mathrm{level}\left(
            \mathcal{M}_{\ell_0,\ldots,\ell_\kappa}
        \right)=\max(\ell_0,\ldots,\ell_\kappa)+1
    }
    \and
    \inferrule[Level-TyInterp]{
        \Gamma\vdash t\Rightarrow\mathcal{M}_{
            \ell_0,\ldots,\ell_\kappa}\\
        \Gamma\vdash\Theta\\
        \Gamma;\Theta\vdash\mathrm{level}(T)=j\\
        j\leqslant\kappa
    }{\Gamma\vdash\mathrm{level}\left(
        T_\Delta^t\right) = \ell_j}
\end{mathpar}
\caption{
    Context well-formedness \& level placement rules}
\Description{}
\end{figure}

\begin{figure}
\centering
\begin{mathpar}
    \inferrule[Infer-Uni]{\Gamma\;\mathrm{wf}}{
        \Gamma\vdash\mathtt{U}_\ell\Rightarrow
            \mathcal{U}_{\ell+1}}
    \and
    \inferrule[Infer-Nat]{\Gamma\;\mathrm{wf}}{
        \Gamma\vdash\mathtt{Nat}_\ell\Rightarrow
            \mathcal{U}_\ell}
    \and
    \inferrule[Check-Zero]{\Gamma\;\mathrm{wf}}
        {\Gamma\vdash 0\Leftarrow\mathrm{Nat}_\ell}
    \and
    \inferrule[Check-Succ]
        {\Gamma\vdash t\Leftarrow\mathrm{Nat}_\ell}
        {\Gamma\vdash\mathrm{S}\,t\Leftarrow
            \mathrm{Nat}_\ell}
    \and
    \inferrule[Infer-NatInd]{
        \Gamma\vdash t\Rightarrow\mathrm{Nat}_\ell\\
        \Gamma,x:\mathrm{Nat}_\ell\vdash
            \mathrm{level}(T)=\_\\
        \Gamma\vdash s_0\Leftarrow T[x\coloneqq 0]\\
        \Gamma,y:\mathrm{Nat}_\ell,z:T[x\coloneqq y]
            \vdash s_S\Leftarrow
                T[x\coloneqq\mathrm{S}\,y]
    }{\Gamma\vdash
        \mathrm{natind}_{x.T}\,t\,\{s_0,yz.s_S\}
        \Rightarrow T[x\coloneqq t]}
    \and
    \inferrule[Infer-Sigma]{
        \Gamma\vdash s\Rightarrow\mathcal{U}_\kappa\\
        \Gamma,x:\mathrm{El}_\kappa(s)\vdash t
            \Rightarrow\mathcal{U}_\ell
    }{\Gamma\vdash\mathtt{Sigma}(x:s)t\Rightarrow
        \mathcal{U}_{\max(\kappa,\ell)}}
    \and
    \inferrule[Check-Pair]{
        \Gamma\vdash s\Leftarrow S\\
        \Gamma\vdash t\Leftarrow T[x\coloneqq s]
    }{\Gamma\vdash(s,t)\Leftarrow(x:S)\times T}
    \and
    \inferrule[Infer-$\mathrm{Proj}_1$]
        {\Gamma\vdash t\Rightarrow(x:S)\times T}
        {\Gamma\vdash \pi_1\,t\Rightarrow S}
    \and
    \inferrule[Infer-$\mathrm{Proj}_2$]
        {\Gamma\vdash t\Rightarrow(x:S)\times T}
        {\Gamma\vdash \pi_2\,t\Rightarrow
            T[x\coloneqq \pi_1\,t]}
    \and
    \inferrule[Infer-Pi]{
        \Gamma\vdash s\Rightarrow\mathcal{U}_\kappa\\
        \Gamma,x:\mathrm{El}_\kappa(s)\vdash t
            \Rightarrow\mathcal{U}_\ell
    }{\Gamma\vdash\mathtt{Pi}(x:s)t\Rightarrow
        \mathcal{U}_{\max(\kappa,\ell)}}
    \and
    \inferrule[Check-Abs]
        {\Gamma,x:S\vdash t\Leftarrow T}
        {\Gamma\vdash\lambda x.\,t\Leftarrow (x:S)\to T}
    \and
    \inferrule[Infer-App]{
        \Gamma\vdash t\Rightarrow (x:S)\to T\\
        \Gamma\vdash s\Leftarrow S
    }{\Gamma\vdash t\,s\Rightarrow T[x\coloneqq s]}
    \and
    \inferrule[Infer-Id]{
        \Gamma\vdash t\Rightarrow\mathcal{U}_\ell\\
        \Gamma\vdash t_1,t_2\Leftarrow
            \mathrm{El}_\ell(t)
    }{\Gamma\vdash\mathtt{Id}\,t\,t_1\,t_2\Rightarrow
        \mathcal{U}_\ell}
    \and
    \inferrule[Check-Refl]
        {\Gamma\vdash t_1\equiv t_2:T}
        {\Gamma\vdash\mathrm{refl}\Leftarrow t_1=_T t_2}
    \and
    \inferrule[Infer-J]{
        \Gamma\vdash t\Rightarrow t_1=_T t_2\\
        \Gamma,x:T,y:t_1=_T x\vdash\mathrm{level}(S)=
            \_\\
        \Gamma\vdash s\Leftarrow
            S[x\coloneqq t_1,y\coloneqq\mathrm{refl}]
    }{\Gamma\vdash\mathrm{J}_{xy.S}\,t\,s\Rightarrow
        S[x\coloneqq t_2,y\coloneqq t]}
    \and
    \inferrule[Infer-Var]
        {\text{\faLock}\not\in\Delta}
        {\Gamma,x:T,\Delta\vdash x\Rightarrow T}
    \and
    \inferrule[Infer-Box]
        {\Gamma.\,\text{\faLock}\vdash t\Rightarrow
            \mathcal{U}_\ell}
        {\Gamma\vdash\mathtt{Box}\,t\Rightarrow
            \mathcal{U}_\ell}
    \and
    \inferrule[Check-Lock]
        {\Gamma.\,\text{\faLock}\vdash t\Leftarrow T}
        {\Gamma\vdash[t]_\text{\faLock}\Leftarrow\Box T}
    \and
    \inferrule[Infer-Unlock]
        {\Gamma^\text{\faUnlock}\vdash t\Rightarrow
            \Box T}
        {\Gamma\vdash[t]_\text{\faUnlock}\Rightarrow T}
    \and
    \inferrule[Infer-Interp]{
        \Gamma\vdash t\Rightarrow\Box T\\
        \Gamma.\,\text{\faLock}\vdash
            \mathrm{level}(T)=\kappa
    }{\Gamma\vdash\llbracket t\rrbracket^{
        \ell_0,\ldots,\ell_\kappa}\Rightarrow
        \left(\mathtt{m}:\mathcal{M}_{
            \ell_0,\ldots,\ell_\kappa}\right)\to
        T^\mathtt{m}_\cdot
    }
    \and
    \inferrule[Infer-Interm-Interp]{
        \Gamma\vdash s\Rightarrow\mathcal{M}_{
            \ell_0,\ldots,\ell_\kappa}\\
        \Gamma.\,\text{\faLock},\Delta\vdash t
            \Rightarrow T
    }{\Gamma\vdash t^s_\Delta\Rightarrow T^s_\Delta}
    \\
    \inferrule[Infer-Ascr]{
        \Gamma\vdash\mathrm{level}(T)=\_\\
        \Gamma\vdash t\Leftarrow T
    }{\Gamma\vdash t:T\Rightarrow T}
    \and
    \inferrule[Check-Inferred]{
        \Gamma\vdash t\Rightarrow S\\
        \Gamma\vdash S\equiv T
    }{\Gamma\vdash t\Leftarrow T}
\end{mathpar}
\caption{Bidirectional term typing rules}
\Description{}
\end{figure}

\subsection{Metatheoretical properties of MLTT-R}
\label{MLTT-R-properties}

Once again, we only try to make sure of the most
essential properties: that computations do not diverge,
that types during computations are preserved, all
computations terminate and that resulting type system
is a conservative extension of MLTT.

\begin{theorem}[Church-Rosser theorem for STLC-R]
    Reflexive transitive closure of $(\to)$ has
    Church-Rosser property.
\end{theorem}
\begin{proof}
    Classic proof via parallel reduction relation
    \cite{CRproof}.
\end{proof}

\begin{theorem}[Type preservation for MLTT-R]
\end{theorem}
\begin{proof}
    By structural induction on terms and $(\to)$.
\end{proof}

\begin{theorem}[Strong normalization for MLTT-R]
    Closed terms typable in MLTT-R are strongly
    normalizing with respect to $(\to)$.
\end{theorem}
\begin{proof}
    Tait's method of logical relations for MLTT with
    modifications similar to ones performed in a proof
    of strong normalization for STLC-R should suffice. 
    However, we're not completely sure.
\end{proof}

\begin{theorem}[Conservativity over MLTT]
    For MLTT context $\Gamma$, MLTT type $T$ and MLTT-R
    term $t$, if $\Gamma\vdash t\Leftarrow T$, then
    $\Gamma\vdash_\mathrm{MLTT}\mathrm{nf}(t):T$.
\end{theorem}
\begin{proof}
    Thanks to type preservation, normal form
    $\mathrm{nf}(t)$ is typable in MLTT-R. TODO TODO
    TODO. Therefore $\mathrm{nf}(t)$ is an MLTT term and    is typable in MLTT.
\end{proof}

\subsection{Useful reinterpretations inside MLTT-R}
\label{MLTT-R-applications}

Here we list some of the language features which
reinterpretation enables without any additional language
support.

\subsubsection{Crude, but effective stratification}
\label{MLTT-R-displacement}

Due to McBride \cite{McBride:Crude2011,
McBride:Crude2012}, a technique of crude, but effective
stratification works as follows: instead of providing
complete support for level polymorphism, let's consider
a \textit{displacement operator} which, given a term $t$
using universes $U_i,U_j,U_k,\ldots$, produces a term
which is equal to $t$ except for universes which are
replaced with $U_{i+1},U_{j+1},U_{k+1},\ldots$. The
displaced term is also well-typed, however it is defined
<<one level higher>>. Notice that displacement operator
is \textit{definable} in MLTT-R as
$\_\,@\,\mathcal{U}_\mathrm{McBride}$ where

\[\mathcal{U}_\mathrm{McBride} = \langle
  \Lambda l.U_{suc\,l}, \ldots
\rangle\]

\subsubsection{Displayed type theory}
\label{MLTT-R-displayedtt}

\subsubsection{Internal external parametricity}
\label{MLTT-R-parametricity}

Demonstrated by Reynolds \cite{reynolds1983types},
parametricity is a property of a type system that, in
simple words, <<polymorphic function cannot meaningfully
depend on generic inputs>>. Thus any polymorphic
function of type $\sigma$ satisfies a certain <<free
theorem>> $\llbracket \sigma \rrbracket$ defined via
relational interpretation of a type.

In recent works \cite{Bernardy:Parametricity2010,
Altenkirch:Parametricity2024}, a streamlined definition
is given for parametricity interpretation of
dependently-typed systems dubbed <<external
parametricity>>. Once again, this relational
interpretation is definable in our system
\textit{internally} with the help of a relational model
$\mathcal{U}_\mathrm{para}$. Since interpreted functions
do not depend on local context, we sidestep issues with
relational context highlighted by Altenkirch and others
\cite{Altenkirch:Parametricity2024}, thus result of the
interpretation is a proper proof of a free theorem.

\subsubsection{Synthetic mathematics}
\label{MLTT-R-synthetic}

As thoroughly discussed in a recent paper by Nawrocki
and others \cite{Nawrocki:Lean2026}, thanks to MLTT
being an internal language of certain
$(\infty,1)$-categories, we can do \textit{synthetic
mathematics} inside it: prove a theorem inside MLTT
about types; interpret it in a model of MLTT and yield a
proof about objects in that model. It is easy to see
that reinterpretation operator we have in MLTT-R is
sufficient to do synthetic mathematics in it and
immediately apply results interpreted in a suitable
model.
